\section{Objetivos}

\subsection{Objetivo general}
Diseñar e implementar una maqueta de un Car Wash Inteligente basada en una red de sensores y actuadores controlada por tres microcontroladores ATmega328P y un ESP32, capaz de automatizar el recorrido de un automóvil de juguete, monitorear variables del proceso, visualizar información localmente y transmitir las mediciones hacia Adafruit IO mediante WiFi.

\subsection{Objetivos específicos}
\begin{itemize}[leftmargin=*]
    \item Implementar una red I2C entre un ATmega328P Maestro, dos ATmega328P Slaves y al menos un sensor con interfaz I2C.
    \item Integrar tres sensores de distinto tipo con funciones específicas dentro del proceso: detección de presencia, monitoreo del nivel de agua y detección de la posición final del automóvil.
    \item Controlar tres tipos de motores requeridos por el proyecto: un motor DC, un motor paso a paso y servomotores empleados en mecanismos del proceso.
    \item Mostrar desde el microcontrolador Maestro información relevante del sistema en una pantalla LCD de 16x2.
    \item Coordinar el funcionamiento simultáneo de sensores, actuadores y comunicaciones para ejecutar automáticamente la secuencia del lavado.
    \item Transferir las variables de medición al ESP32 mediante el protocolo de comunicación definido en la implementación final y enviarlas en tiempo real a Adafruit IO mediante WiFi.
    \item Diseñar y construir una maqueta en MDF que permita montar los sensores y actuadores y guiar el desplazamiento del automóvil durante el proceso.
\end{itemize}
